\section{Encode pictures that another program can decode}
The previous chapter defined both sides of a small lossless format. Here, the
output must follow a format implemented elsewhere. A file extension and a familiar
start code are insufficient: another decoder must recover the intended pictures,
in the intended order, from your emitted syntax. Keep that independent decoder
outside the MARS implementation and outside your own reconstruction code.

Build an encoder for a restricted baseline H.263 stream. The input is a short,
script-generated QCIF sequence of raw samples. The target output uses I-pictures
and P-pictures, with optional coding modes disabled. An I-picture establishes a
reference without an earlier picture; a P-picture can describe predictions and
residuals relative to the preceding reconstructed reference. This profile uses
neither the earlier chapter's XOR commands nor a private teaching container.

The new difficulty is controlled loss. A quantized block usually reconstructs to
values different from its input. Measure that error, but keep its meaning separate
from malformed syntax, incorrect memory bounds or a wrong prediction reference.
A successfully decoded file can still contain poor pictures. Conversely, two
valid encoders can choose different predictions and produce different bytes for
similar pictures. Final feedback must examine syntax, reconstruction and quality,
rather than require one reference bitstream.

Begin with the input layout and an exact block-comparison kernel. The kernel is
small enough to inspect completely, yet it will be called repeatedly during motion
search. Then connect prediction, transform, quantization, local reconstruction and
bit writing. Students are expected to know MSA instructions already; the questions
here concern which values belong in their lanes and which bounds make their use
correct.
